% Options for packages loaded elsewhere
% Options for packages loaded elsewhere
\PassOptionsToPackage{unicode}{hyperref}
\PassOptionsToPackage{hyphens}{url}
\PassOptionsToPackage{dvipsnames,svgnames,x11names}{xcolor}
%
\documentclass[
  11pt,
  a4paper,
]{article}
\usepackage{xcolor}
\usepackage[a4paper,textheight=24cm,textwidth=15.5cm]{geometry}
\usepackage{amsmath,amssymb}
\setcounter{secnumdepth}{5}
\usepackage{iftex}
\ifPDFTeX
  \usepackage[T1]{fontenc}
  \usepackage[utf8]{inputenc}
  \usepackage{textcomp} % provide euro and other symbols
\else % if luatex or xetex
  \usepackage{unicode-math} % this also loads fontspec
  \defaultfontfeatures{Scale=MatchLowercase}
  \defaultfontfeatures[\rmfamily]{Ligatures=TeX,Scale=1}
\fi
\usepackage[]{libertinus}
\ifPDFTeX\else
  % xetex/luatex font selection
  \setmonofont[Scale=0.92]{Latin Modern Mono}
\fi
% Use upquote if available, for straight quotes in verbatim environments
\IfFileExists{upquote.sty}{\usepackage{upquote}}{}
\IfFileExists{microtype.sty}{% use microtype if available
  \usepackage[]{microtype}
  \UseMicrotypeSet[protrusion]{basicmath} % disable protrusion for tt fonts
}{}
\makeatletter
\@ifundefined{KOMAClassName}{% if non-KOMA class
  \IfFileExists{parskip.sty}{%
    \usepackage{parskip}
  }{% else
    \setlength{\parindent}{0pt}
    \setlength{\parskip}{6pt plus 2pt minus 1pt}}
}{% if KOMA class
  \KOMAoptions{parskip=half}}
\makeatother
% Make \paragraph and \subparagraph free-standing
\makeatletter
\ifx\paragraph\undefined\else
  \let\oldparagraph\paragraph
  \renewcommand{\paragraph}{
    \@ifstar
      \xxxParagraphStar
      \xxxParagraphNoStar
  }
  \newcommand{\xxxParagraphStar}[1]{\oldparagraph*{#1}\mbox{}}
  \newcommand{\xxxParagraphNoStar}[1]{\oldparagraph{#1}\mbox{}}
\fi
\ifx\subparagraph\undefined\else
  \let\oldsubparagraph\subparagraph
  \renewcommand{\subparagraph}{
    \@ifstar
      \xxxSubParagraphStar
      \xxxSubParagraphNoStar
  }
  \newcommand{\xxxSubParagraphStar}[1]{\oldsubparagraph*{#1}\mbox{}}
  \newcommand{\xxxSubParagraphNoStar}[1]{\oldsubparagraph{#1}\mbox{}}
\fi
\makeatother


\usepackage{longtable,booktabs,array}
\newcounter{none} % for unnumbered tables
\usepackage{calc} % for calculating minipage widths
% Correct order of tables after \paragraph or \subparagraph
\usepackage{etoolbox}
\makeatletter
\patchcmd\longtable{\par}{\if@noskipsec\mbox{}\fi\par}{}{}
\makeatother
% Allow footnotes in longtable head/foot
\IfFileExists{footnotehyper.sty}{\usepackage{footnotehyper}}{\usepackage{footnote}}
\makesavenoteenv{longtable}
\usepackage{graphicx}
\makeatletter
\newsavebox\pandoc@box
\newcommand*\pandocbounded[1]{% scales image to fit in text height/width
  \sbox\pandoc@box{#1}%
  \Gscale@div\@tempa{\textheight}{\dimexpr\ht\pandoc@box+\dp\pandoc@box\relax}%
  \Gscale@div\@tempb{\linewidth}{\wd\pandoc@box}%
  \ifdim\@tempb\p@<\@tempa\p@\let\@tempa\@tempb\fi% select the smaller of both
  \ifdim\@tempa\p@<\p@\scalebox{\@tempa}{\usebox\pandoc@box}%
  \else\usebox{\pandoc@box}%
  \fi%
}
% Set default figure placement to htbp
\def\fps@figure{htbp}
\makeatother


% definitions for citeproc citations
\NewDocumentCommand\citeproctext{}{}
\NewDocumentCommand\citeproc{mm}{%
  \begingroup\def\citeproctext{#2}\cite{#1}\endgroup}
\makeatletter
 % allow citations to break across lines
 \let\@cite@ofmt\@firstofone
 % avoid brackets around text for \cite:
 \def\@biblabel#1{}
 \def\@cite#1#2{{#1\if@tempswa , #2\fi}}
\makeatother
\newlength{\cslhangindent}
\setlength{\cslhangindent}{1.5em}
\newlength{\csllabelwidth}
\setlength{\csllabelwidth}{3em}
\newenvironment{CSLReferences}[2] % #1 hanging-indent, #2 entry-spacing
 {\begin{list}{}{%
  \setlength{\itemindent}{0pt}
  \setlength{\leftmargin}{0pt}
  \setlength{\parsep}{0pt}
  % turn on hanging indent if param 1 is 1
  \ifodd #1
   \setlength{\leftmargin}{\cslhangindent}
   \setlength{\itemindent}{-1\cslhangindent}
  \fi
  % set entry spacing
  \setlength{\itemsep}{#2\baselineskip}}}
 {\end{list}}
\usepackage{calc}
\newcommand{\CSLBlock}[1]{\hfill\break\parbox[t]{\linewidth}{\strut\ignorespaces#1\strut}}
\newcommand{\CSLLeftMargin}[1]{\parbox[t]{\csllabelwidth}{\strut#1\strut}}
\newcommand{\CSLRightInline}[1]{\parbox[t]{\linewidth - \csllabelwidth}{\strut#1\strut}}
\newcommand{\CSLIndent}[1]{\hspace{\cslhangindent}#1}



\setlength{\emergencystretch}{3em} % prevent overfull lines

\providecommand{\tightlist}{%
  \setlength{\itemsep}{0pt}\setlength{\parskip}{0pt}}

\newcommand{\specialcell}[2][c]{%
	\begin{tabular}[#1]{@{}c@{}}#2\end{tabular}}
\newcommand\indica{\,\rlap{{\rm \Large 1}}\kern.30em \mbox{\rm \Large l}}

 


\usepackage{amsmath}
\usepackage{tikz}
\usepackage{xcolor}
\usepackage{tabularx}
\usepackage{orcidlink}
\usepackage{lineno}
\usepackage{unicode-math}
\setmathfont{Libertinus Math}
\usepackage{draftwatermark}
\usepackage{adjustbox}
\usepackage{float}


\newcommand\C{{\mathbb C}}
\newcommand\E{{\mathbb E}}
\newcommand\F{{\mathbb F}}
\newcommand\M{{\mathbb M}} % un espace metrique
\newcommand\N{{\mathbb N}}
\newcommand\R{{\mathbb R}}

\newtheorem{theorem}{Theorem}    
\newtheorem{acknowledgement}[theorem]{Acknowledgment}    
\newtheorem{algorithm}[theorem]{Algorithm}    
\newtheorem{axiom}[theorem]{Axiom}    
\newtheorem{case}[theorem]{Case}    
\newtheorem{claim}[theorem]{Claim}    
\newtheorem{conclusion}[theorem]{Conclusion}    
\newtheorem{condition}[theorem]{Condition}    
\newtheorem{conjecture}[theorem]{Conjecture}    
%\newtheorem{corollary}[theorem]{Corollary} 
\newtheorem{corollary}{Corollary}   
\newtheorem{criterion}[theorem]{Criterion}    
%\newtheorem{definition}[theorem]{Definition} 
\newtheorem{definition}{Definition}  
\newtheorem{example}[theorem]{Example}    
\newtheorem{exercise}[theorem]{Exercise}    
\newtheorem{lemma}{Lemma}     
\newtheorem{notation}[theorem]{Notation}    
\newtheorem{problem}[theorem]{Problem}    
\newtheorem{proposition}{Proposition}
%\newtheorem{lemma}{Lemma}    
\newtheorem{remark}[theorem]{Remark}    
\newtheorem{solution}[theorem]{Solution}    
\newtheorem{summary}[theorem]{Summary}   


% Tables converted from `longtable` to `table` (see pdf-table-fit.lua, used
% to shrink-to-fit overly wide tables) must stay in place like `longtable`
% does, instead of floating to the top of the page/document.
\floatplacement{table}{H}

\makeatletter
\@ifpackageloaded{caption}{}{\usepackage{caption}}
\AtBeginDocument{%
\ifdefined\contentsname
  \renewcommand*\contentsname{Table of contents}
\else
  \newcommand\contentsname{Table of contents}
\fi
\ifdefined\listfigurename
  \renewcommand*\listfigurename{List of Figures}
\else
  \newcommand\listfigurename{List of Figures}
\fi
\ifdefined\listtablename
  \renewcommand*\listtablename{List of Tables}
\else
  \newcommand\listtablename{List of Tables}
\fi
\ifdefined\figurename
  \renewcommand*\figurename{Figure}
\else
  \newcommand\figurename{Figure}
\fi
\ifdefined\tablename
  \renewcommand*\tablename{Table}
\else
  \newcommand\tablename{Table}
\fi
}
\@ifpackageloaded{float}{}{\usepackage{float}}
\floatstyle{ruled}
\@ifundefined{c@chapter}{\newfloat{codelisting}{h}{lop}}{\newfloat{codelisting}{h}{lop}[chapter]}
\floatname{codelisting}{Listing}
\newcommand*\listoflistings{\listof{codelisting}{List of Listings}}
\makeatother
\makeatletter
\makeatother
\makeatletter
\@ifpackageloaded{caption}{}{\usepackage{caption}}
\@ifpackageloaded{subcaption}{}{\usepackage{subcaption}}
\makeatother
\makeatletter
\@ifpackageloaded{algorithm}{}{\usepackage{algorithm}}
\makeatother
\makeatletter
\@ifpackageloaded{algpseudocode}{}{\usepackage{algpseudocode}}
\makeatother
\makeatletter
\@ifpackageloaded{caption}{}{\usepackage{caption}}
\makeatother
    \makeatletter
    \newcommand\fs@nocaption{
      \def\@fs@cfont{\bfseries}
      \let\@fs@capt\floatc@plain
      \def\@fs@pre{}%
      \def\@fs@post{\kern2pt\hrule}%
      \def\@fs@mid{\hrule\kern2pt}%
      \let\@fs@iftopcapt\iftrue}
    \makeatother
  
\captionsetup[algorithm]{justification=raggedright}
\usepackage{bookmark}
\IfFileExists{xurl.sty}{\usepackage{xurl}}{} % add URL line breaks if available
\urlstyle{same}
% fallback for those not using the hyperref driver hyperxmp:
\makeatletter
\@ifundefined{xmpquote}{\newcommand{\xmpquote}[1]{#1}}{}
\makeatother
\hypersetup{
  pdftitle={Median search: dichotomy and gradient descent},
  pdfauthor={Elie Youndje},
  pdfkeywords={\xmpquote{Mean}, \xmpquote{median}, \xmpquote{order
statistics}, \xmpquote{sort}, \xmpquote{sum of absolute deviations}},
  colorlinks=true,
  linkcolor={blue},
  filecolor={Maroon},
  citecolor={Blue},
  urlcolor={Blue},
  pdfcreator={LaTeX via pandoc}}


\title{Median search: dichotomy and gradient descent}
\usepackage{etoolbox}
\makeatletter
\providecommand{\subtitle}[1]{% add subtitle to \maketitle
  \apptocmd{\@title}{\par {\large #1 \par}}{}{}
}
\makeatother
\subtitle{Dichotomy and gradient algorithms}
\author{Elie Youndje}
\date{2026-11-11}
\begin{document}
\definecolor{computo-blue}{HTML}{034E79}

% Neutral header for drafts/submissions: no Computo branding, so the PDF can
% travel to HAL, PCI, arXiv, double-blind reviewers, etc. without implying
% any relationship to the journal beyond a discreet mention.
\begin{center}
  {\Large\bfseries Median search: dichotomy and gradient descent}
  \\[.3cm]{\large Dichotomy and gradient algorithms}
\end{center}

\bigskip
\begin{center}
      \textbf{Elie
Youndje}\textsuperscript{1}~\orcidlink{0000-0001-8546-4530}{\renewcommand{\thefootnote}{\fnsymbol{footnote}}\footnote{Corresponding author: \href{mailto:elie.youndje@univ-rouen.fr}{elie.youndje@univ-rouen.fr}}}\quad
  
  \smallskip
  {\small
      \textsuperscript{1}Mathematics, Univ Rouen Normandie, CNRS,
Normandie Univ, LMRS UMR 6085\quad
    }

  \bigskip

  Date published: 2026-11-11 \quad Last modified: 2026-11-11
  \end{center}

\bigskip
\begin{abstract}
This paper shows how to compute the median of a sample by minimizing the
sum of absolute deviations function. The minimization is done using the
left hand gradient of the sum of absolute deviations. The point where
the left hand derivative changes sign is searched using dichotomy or
gradient descent. Simulations conducted in the article reveal that these
approaches are very promising. In particular, for each sample size
considered, there is one among our methods that outperforms the matlab
build-in function in terms of runtime.
\end{abstract}

\noindent%
{\it Keywords:} Mean, median, order statistics, sort, sum of absolute
deviations
\vfill

\DraftwatermarkOptions{stamp=true, colorspec=0.9, text={\rm submitted}}
\linenumbers

\floatname{algorithm}{Algorithm}

\renewcommand*\contentsname{Contents}
{
\setcounter{tocdepth}{3}
\tableofcontents
}

\section{Introduction}\label{introduction}

Let \(X\) be a random variable, \(X_1,\ldots,X_n\) be an i.i.d. sample
from \(X.\) Let \(X_{(1)},\ldots,X_{(n)}\) be the order statistics
associated with \(X_1,\ldots,X_n.\) As we learnt in school the
(empirical) mean (\(\bar{\mu}_0\)) and median (\(\bar{\mu}_1\)) of \(X\)
are defined respectively by \[\bar{\mu}_0=\frac{1}{n}\sum_{i=1}^n X_i,\]
\[\bar{\mu}_1   = \begin{cases}
	X_{(m+1)}  & \text{if}\  n=2m+1, \\
	\frac{X_{(m)}+X_{(m+1)}}{2} & \text{if}\ n=2m.
\end{cases}\] Both quantities \(\bar{\mu}_0\) and \(\bar{\mu}_1\)
are used in practice as estimates of the ``center" of the distribution
of \(X.\) Let \(F\) be the cumulative distribution function of \(X,\)
the theoretical mean of \(X\) defined by \[\mu_0=\int x \mathrm{d}F(x)\]
does not always exist. To the contrary the theoretical median of \(X\)
given by \[\mu_1=\inf\{ x:\ \  F(x)\ge \frac{1}{2}\}\] always exists.
Besides this existence issue, medians are very important in practice
because they are resistant to outliers. This means that they are not
heavily affected by a very high or a very low value of the data. Median
computation is used in many fields. For seismic data processing we can
cite (Stewart 1985). In image processing, median filters are preferred
to mean filters because of their edge enhancement properties.
Applications involving median computation can be found in (Huang et al.
1979), (Nodes and Gallagher 1982), (Azzini et al. 2023) and (Yin et al.
1996). As mention in the literature (see for instance (Tibshirani
2009)), median is more appropriate than mean to measure the center of a
population. However the appealing qualities of the median come with a
high price: computing the median takes more time than computing the
mean. For this reason, many methods for computing median have been
devised. The methods based on the''divide and conquer" principal are
developed in (Hoare 1961), (Tibshirani 2009) and (Azzini et al. 2023).
These methods are very interesting and significantly faster than the
``sort then select" method implemented in most computer languages.
However they do not always give us what we expect to get. Indeed, when
\(n=2m\), they return \(X_{(m)}\) or \(X_{(m+1)}\) and none of these
quantities is the median given in textbooks. In (Satman 2022) is
introduced a novel iterative method to approximate the median. Although
this method appears to be interesting, it is not shown in the paper, how
to choose the parameter used in its construction. Another class of
iterative methods that can be used to compute the median of a data set
is the subgradient. The sugradient method is a conventional means to
minimize a nonsmooth function, therefore it can be used to compute the
median by minimizing the sum of absolute deviations. Good expositions of
the subgradient algorithm can be found in (Tibshirani 2017) and (Ma
2019). The class of subgradient methods shares some characteristics of
the methods introduce in this paper:

\begin{itemize}
	\item[--] They build sequences ``converging" to the median.
	\item[--] The sequences are build using the left hand side (one-sided) derivative of the sum of absolute deviations.
\end{itemize}
The approaches of this paper and the class of subgradient methods also have some ingredients or properties that are very
different. For example:
\begin{itemize}
	\item[--] The subgradient class  builds  a sequence $(x_n,n\in \N)$ such that
	\[ \lim x_n=\bar{\mu}_1 \]
	where $\bar{\mu}_1$ is the sample median. While the class of methods developed in this article builds  two (adjacent)
	sequences $(a_n,n\in \N), \quad (b_n,n\in \N) $ such that 
	\[ a_n\leq \bar{\mu}_1 \leq b_n \quad \text{and}\quad \lim( b_n-a_n)=0.\]
	
	\item[--] The output of our approaches is the \textit{exact} median $\bar{\mu}_1,$ whereas the output of the subgradient class
	is an \textit{approximate} value of $\bar{\mu}_1.$
\end{itemize}
The aim of this
article is to introduce  methods which return the median as they  teach it in school and at the same time, on average,
outperforms the conventional ``sorting" method for ``large'' sample size.

The next section contains  theoretical foundations of the  methods introduced in this paper. The common principles of our algorithms
are presented in Section~3. The dichotomic method is outlined in Section 4 and the last section of the paper contains the gradient descent
method.






\section{Theoretical preliminaries }\label{sec:2}
Let $x\in\R$, the sum of absolute deviation of the sample $X_1,\ldots,X_n$ from $x$ is given
by
\[
\varphi(x)=\sum_{i=1}^{n}\lvert X_i-x\rvert.
\]
It is well-known that the median of $X_1,\ldots,X_n$ is a minimizer of the  function $\varphi$, in
others words
\[
\varphi( \bar{\mu}_1) =\inf_{x} \varphi(x).
\]
Lemma 1 below is fundamental, it gives us useful information about $\varphi.$ The proof of Lemma 1 is in the supplementary material. In order to state this result
we introduce the following notations. Let $Y_1<Y_2<\ldots<Y_m$  be such that
\[
\{X_1,\ldots,X_n\}=\{Y_1,\ldots,Y_m\}.
\]
Set 
\begin{align*}
	S = &\{Y_1,\ldots,Y_m\}, \\
	I_0= & (-\infty,Y_1), \quad I_1=[Y_1,Y_2),\ldots, I_{m-1}=[Y_{m-1},Y_{m}),\quad I_m=[Y_m,+\infty).
\end{align*}
%%%%%      
\begin{lemma} {\ }
	\begin{itemize}
		\item[(i)] For $k\in \{0,\ldots,m\},$  $\exists \alpha_k, \beta_k\in\R$ such that 
		\[
		\varphi(x)=\alpha_k x+ \beta_k, \quad \forall x\in I_k.\quad \quad\quad\quad
		\]	
		\item[(ii)]  For $j\in\{1,\ldots,m\}$, $\exists q_j\in\N^*$ such that
		\[
		\alpha_j- \alpha_{j-1}=2 q_j.
		\]      
		\item[(iii)] $\quad \quad-n=\alpha_0 <\alpha_1<\ldots <\alpha_{m-1}<\alpha_m=n.	$
	\end{itemize}	     
\end{lemma} 





We are now going to introduce the definitions and notations, we will use to state the next theorem.
The right hand derivative of $\varphi$ at $a\in \R$ is defined by
\[
\varphi'_+(a)=\lim_{\substack{h\to 0 \\ h>0}} \frac{\varphi(a+h)-\varphi(a)}{h}.
\]
Similarly the left hand derivative is given by
\[
\varphi'_-(a)=\lim_{\substack{h\to 0 \\ h<0}} \frac{\varphi(a+h)-\varphi(a)}{h}.
\]
Let $a_0\in\R,$ $a_0$ is called
\begin{itemize}
	\item  [--] a \textbf{mininum} of $\varphi$ if
	\[
	\forall x\in\R,\  \varphi(a_0)\leq \varphi(x),
	\]
	\item  [--] an \textbf{absolute mininum (am)} of $\varphi$ if
	\[
	\forall x\in\R\setminus\{a_0\},\  \varphi(a_0)< \varphi(x).
	\]        
\end{itemize}
%%%%%%%%%%%%
In the sequel we dote by $\mathscr{M}$ the set of minimum of $\varphi$ \i.e.
\[
\mathscr{M}=\{a\in\R,\ a \ \text{minimun of} \ \varphi\}.  
\]
We make use of the following definitions below. 
%%%%%%%%%%%%%
\begin{itemize}
	\item  [--] A point $a_1$ of $\R$ is called the   \textbf{least mininum (lm)} of $\varphi$ if
	\[
	a_1=\min\left(\mathscr{M}\right).
	\]
	\item  [--] A point $a_2$ of $\R$ is called the   \textbf{greatest mininum (gm)} of $\varphi$ if
	\[
	a_2=\max\left(\mathscr{M}\right).
	\]
\end{itemize}
We have the following foundation  result whose proof is in the supplementary material.  

\begin{theorem} For any data set $S=\{X_1,\ldots,X_n\},$ one and only one of the following two statements  ((a), (b)) is true.
	\begin{itemize}
		\item[(a)] $\varphi$ has an \textbf{absolute mininum (am)} $a_0$ satisfying:
		\begin{itemize}
			\item[(i)] $a_0\in S,$
			\item[(ii)] $a_0$ is the only point satisfying
			\[
			\varphi'_-(a_0)<0 \quad\text{and}\quad \varphi'_+(a_0)>0.
			\]
		\end{itemize}	  
		\item[(b)] $\varphi$ has a \textbf{lm} $a_1$    and a \textbf{gm} $a_2.$ Furthermore
		in this case we have:
		\begin{itemize}
			\item[(i)] $a_1,\ a_2\in S,$
			\item[(ii)] $\mathscr{M}=[a_1,\ a_2]$
			
			\item[(iii)] $a_1$ is the only point satisfying
			\[
			\varphi'_-(a_1)<0 \quad\text{and}\quad \varphi'_+(a_1)=0,
			\]
			\item[(iv)] $a_2$ is the only point satisfying
			\[
			\varphi'_-(a_2)=0 \quad\text{and}\quad \varphi'_+(a_2)>0.
			\]
		\end{itemize}	 
	\end{itemize}	
	
\end{theorem} 


\begin{corollary}  For a data set $S=\{X_1,\ldots,X_n\},$  with $n=2p+1,$
	$\varphi$ has an \textbf{absolute mininum (am)} $a_0$ satisfying:
	\begin{itemize}
		\item[(i)] $a_0\in S,$
		\item[(ii)] $a_0$ is the only point satisfying
		\[
		\varphi'_-(a_0)<0 \quad\text{and}\quad \varphi'_+(a_0)>0.
		\]
	\end{itemize}	
\end{corollary}





We are now going to introduce algorithmically convenient formulas for the left hand and right hand derivative
of the function $\varphi.$ To state these results we need some notations. For $v\in \R,$ set
\begin{align*}
	L_v= &\{i:\  X_i<v\},\quad R_v=\{i:\  X_i>v\},\\
	l_v= &\#(L_v)\quad \text{and}\quad r_v= \#(R_v).
\end{align*}
We have the following lemma.


\begin{lemma} For $v\in \R,$ we have:
	\[
	\varphi'_-(v)=-n+2l_v \quad\text{and}\quad \varphi'_+(v)=n-2r_v . \]
\end{lemma}



Given $u,v\in \R,$ the next result gives the relationships between $l_u$ and $l_v$ and between $r_u$ and $r_v$.
These relationships will help in our algorithms, to speed   up the computation of $ \varphi'_-(v)$  when    $\varphi'_-(u)$
is known. To state this result, we introduce the following notations. For $u,v \in \R,$ $u<v$ we set
\begin{align*}
	E_u= &\{i:\ X_i=v\},\quad e_u=\#(E_u),\\
	S_{u,v}=&\{i:\ u< X_i<v\},\quad s_{u,v} =\#(S_{u,v}).
\end{align*}
If $l_v,\ r_v,\ e_v, \ e_u,$ and $s_{u,v}$ are known, the lemma below shows that we can deduce
$l_u$ and $r_u.$


\begin{lemma} For $u<v$ we have:
	\begin{align*}
		(i)\  \, l_u=& l_v-e_u-s_{u,v},\\
		(ii)\   r_u=& r_v+e_v+s_{u,v}.
	\end{align*}
\end{lemma}
The proofs of Lemma 2 and Lemma 3 are in the supplementary material.


\section{Basic principles of our methods }\label{sec:3}

In this section, we introduce notations and properties common to both dichotomy and gradient descent.
Let $f:\R\rightarrow \R$ be a convex and differentiable function. It is well known that a minimizer of
$f$ is a solution of the equation $f'(x)=0.$ When it is not possible to solve $f'(x)=0$ analytically,
usually dichotomy or gradient descent is used to build an approximate solution. In this paper, we want
to minimize the convex function $\varphi$, however $\varphi$ is not differentiable. Nonetheless, $\varphi_{-}'(t)$
exists for all $t.$ In this section we show how we can use $\varphi_{-}'(t)$ to compute \textbf{lm}.

Lemma~4 belows gives a useful characterization of \textbf{lm}; its proof is the supplementary material.
\begin{lemma}
	\[
	\mathbf{lm}=\max\{ v:\quad \varphi_{-}'(v)<0
	\}.
	\]	
\end{lemma}




Lemma 4 is very interesting, it can be used to test whether an interval contains \textbf{lm}. Le $a<b$, we have:
\begin{itemize}
	\item if  $\varphi_{-}'(a)<0$ and $\varphi_{-}'(b)\ge 0$ then $\mathbf{lm}\in [a,\ b[;$
	\item if  $\varphi_{-}'(a)\ge 0$ then $\mathbf{lm}\notin [a,\ b];$
	\item if  $\varphi_{-}'(b)< 0$ then $\mathbf{lm}\notin [a,\ b[.$
\end{itemize} 
The function $\varphi_{-}'(v)$ plays important roles in the rest of the paper. To save typing and space we set
\[
\rho(v)=\varphi_{-}'(v).
\]
We now introduce some notations used in the sequel. Let $U$ be a vector in $\R^p$ that is $U=(u_1,\ldots,u_p)$ for
$I=(i_1,\ldots,i_q)\in \{1,2,\ldots,p\}^q \ (q\leq p)$ we set
\[
U(I)=(u_{i_1},u_{i_2},\ldots,u_{i_q}).
\]
Let $e\in\R,$ $K=\{i\ : \ u_i<e\}=\{j_1,\ldots,J_{q_1}\}$ with $j_1<j_2<\ldots<j_{q_1}.$ We define the vector
$(U<e)$ by
\[
(U<e)=U(J)\quad \text{where}\quad J=(j_1,\ldots,J_{q_1}).
\]
Likewise the vector $(U\ge e)$ is defined by 
\[
(U\ge e)=U(J')
\] 
where
\[
J'=(j_1',\ldots,J_{q_2}'),\quad \{j_1',\ldots,J_{q_2}'\}=\{i\ : \ u_i\ge e\}\quad \text{and}\quad j_1'<j_2'<\ldots<j_{q_2}'.
\]        
Let $V=(v_1,\ldots,v_m)\in \R^m,$ in what follows we write $U\subseteq V$ if every component of $U$ is a component
of $V,$ said otherwise $u_i\in \{v_1,\ldots,v_m\}.$
For a vector $V=(v_1,\ldots,v_m)$ we define $\lvert V\rvert$ as
\[
\lvert V\rvert=\mathrm{card}(\{v_i,\ 1\leq i\leq m \}).
\]

The common ideas used in the methods considered in this article consists in building a sequence
$\left(a_i,b_i,c_i,M_i \right)_{i\in \N^*}$ such that:
\begin{itemize}
	\item[(i)] $a_i\leq c_i\leq b_i,$\quad $\mathbf{lm}\in [a_i,\ b_i],$\quad  $[a_{i+1},\ b_{i+1}] \subseteq [a_i,\ b_i] ;$
	\item[(ii)] $M_i\subseteq X=(X_1,\ldots,X_n),$\quad $M_{i+1}\subseteq M_i;$
	\item[(iii)] $\exists m_0\in \N,\quad a_{m_0}=\mathbf{lm}=b_{m_0}.$
\end{itemize}
Below in Step-1, Step-2,  Step-3 we present the construction of the sequence $\left(a_i,b_i,c_i,M_i \right).$


\begin{itemize}
	\item[Step-1]  \textit{construction of $\left(a_1,b_1,c_1,M_1 \right)$ }
	\begin{itemize}
		\item[$\bullet$] $a_1=\min(S);$
		\item[$\bullet$]  $b_1=\max(S);$
		\item[$\bullet$]  $M_1=(X_1,\ldots,X_n);$
		
		\item[$\bullet$] \texttt{if} $(a_1=b_1)$
		\begin{itemize}
			\item[$\circ$] $c_1= a_1;$
		\end{itemize}
		\item[$\bullet$] \texttt{if} $(a_1<b_1)$ 
		\begin{itemize}     
			\item[$\circ$] \texttt{choose} $c_1\in ]a_1,\ b_1[;$	
		\end{itemize}			
	\end{itemize}            
	
	
	
	\item[Step-2]  \textit{construction of $\left(a_2,b_2,c_2,M_2 \right)$ }: 
	\begin{itemize}
		\item[$\bullet$] $L=(M_1<c_1);$
		\item[$\bullet$] $R=(M_1>=c_1);$ 	           
		
		\item[$\bullet$] \texttt{if} $(\rho(c_1)<0)$
		\begin{itemize}
			\item[$\circ$]  \texttt{if} $(L=\emptyset)$
			\begin{itemize}
				\item[$\diamond$]  $a_2=\min(R);$
				\item[$\diamond$] $b_2=b_1;$
			\end{itemize}             
			\item[$\circ$] \texttt{if} $(L\neq\emptyset)$ 
			\begin{itemize} 
				\item[$\diamond$] $a_2=c_1;$
				\item[$\diamond$] $b_2=b_1;$
			\end{itemize}    
			\item[$\circ$] $M_2=R;$
		\end{itemize}     
		\item[$\bullet$] \texttt{if} $(\rho(c_1)\ge 0)$  
		\begin{itemize}			  
			\item[$\circ$] \texttt{if} $(R=\emptyset)$
			\begin{itemize}
				\item[$\diamond$] $b_2=\max(L);$
				\item[$\diamond$] $a_2=a_1;$
			\end{itemize}   
			\item[$\circ$] \texttt{if} $(R\neq\emptyset)$
			\begin{itemize}
				\item[$\diamond$] $b_2=c_1;$
				\item[$\diamond$] $a_2= a_1;$
			\end{itemize}    
			\item[$\circ$] $M_2=L;$
		\end{itemize} 
		\item[$\bullet$] \texttt{if} $(a_2=b_2)$
		\begin{itemize}
			\item[$\circ$] $c_2= a_2;$
		\end{itemize}   
		\item[$\bullet$] \texttt{if} $(a_2<b_2)$  
		\begin{itemize}
			\item[$\circ$] \texttt{choose} $c_2\in ]a_2,\ b_2[;$	
		\end{itemize}
	\end{itemize}
	
	
	\item[Step-3]  \textit{induction step, construction of $\left(a_{k+1},b_{k+1},c_{k+1},M_{k+1} \right)$  }:
	\begin{itemize}
		\item[$\bullet$] $L=(M_k<c_k);$
		\item[$\bullet$] $R=(M_k>=c_k);$ 	           
		
		\item[$\bullet$] \texttt{if} $(\rho(c_k)<0)$
		\begin{itemize}
			\item[$\circ$]  \texttt{if} $(L=\emptyset)$
			\begin{itemize}
				\item[$\diamond$]  $a_{k+1}=\min(R);$
				\item[$\diamond$] $b_{k+1}=b_k;$
			\end{itemize}             
			\item[$\circ$] \texttt{if} $(L\neq\emptyset)$ 
			\begin{itemize} 
				\item[$\diamond$] $a_{k+1}=c_k;$
				\item[$\diamond$] $b_{k+1}=b_k;$
			\end{itemize}    
			\item[$\circ$] $M_{k+1}=R;$
		\end{itemize}     
		\item[$\bullet$] \texttt{if} $(\rho(c_k)\ge 0)$  
		\begin{itemize}			  
			\item[$\circ$] \texttt{if} $(R=\emptyset)$
			\begin{itemize}
				\item[$\diamond$] $b_{k+1}=\max(L);$
				\item[$\diamond$] $a_{k+1}=a_k;$
			\end{itemize}   
			\item[$\circ$] \texttt{if} $(R\neq\emptyset)$
			\begin{itemize}
				\item[$\diamond$] $b_{k+1}=c_k;$
				\item[$\diamond$] $a_{k+1}= a_k;$
			\end{itemize}    
			\item[$\circ$] $M_{k+1}=L;$
		\end{itemize} 
		\item[$\bullet$] \texttt{if} $(a_{k+1}=b_{k+1})$
		\begin{itemize}
			\item[$\circ$] $c_{k+1}= a_{k+1};$
		\end{itemize}   
		\item[$\bullet$] \texttt{if} $(a_{k+1}<b_{k+1})$  
		\begin{itemize}
			\item[$\circ$] \texttt{choose} $c_{k+1}\in ]a_{k+1},\ b_{k+1}[;$	
		\end{itemize}
	\end{itemize}
	
	
\end{itemize} 	
\begin{definition}
	A sequence $\left(\Omega_k=\left(a_k,b_k,c_k,M_k \right),\quad k\in \N^*\right)$ 
	built recursively by  Step-1---Step-3 presented above is called a generic sequence.
\end{definition}      

In Proposition 1--4 below are listed some interesting properties of the sequence  $\left(a_i,b_i,c_i,M_i \right)_{i\in \N^*}.$
The proofs of these results are in the supplementary material.
\begin{proposition} 
	If $\lvert M_k\rvert\geq 2,\  \lvert M_{k+1}\rvert =\lvert M_k\rvert$ and   $\lvert M_{k+2}\rvert=\lvert M_k\rvert$ then
	\[
	\lvert M_{k+3}\rvert < \lvert M_k\rvert.
	\]
\end{proposition}



\begin{proposition} 
	If $\lvert M_k\rvert\geq 2,$  $\exists j>k$ such that:  
	\[
	\lvert M_{j}\rvert < \lvert M_k\rvert.
	\]
\end{proposition}




\begin{proposition} There exists $p_0\in\N$ 
	such that:  
	\[
	\lvert M_{p_0}\rvert =1.
	\]
\end{proposition}



\begin{proposition} There exists $m_0\in\N$ 
	such that:  
	\[
	a_{m_0}=\mathbf{lm}= b_{m_0}.
	\]
\end{proposition}






\section{Dichotomy}\label{sec:3}
This approach is based on Lemma 4. It is very similar to the classical dichotomy used to find the zero of a function.
The main idea here is to use $\varphi_{-}'(v)$ to build a sequence converging to $\mathbf{lm}.$ The dichotomic
sequence is introduce in Definition 2 below.


\begin{definition}
	A sequence $\left(\mathscr{A}_k=\left(a_k,b_k,c_k,M_k \right),\quad k\in \N^*\right)$ is called a dichotomic
	sequence if it is built recursively by  Step-a1  and  Step-a2 presented in the supplementary material.
\end{definition}    


%%%%%%%%%%%%%% 	
It is clear that a dichotomic sequence is a generic sequence, thus we have: $\lim a_k=\mathbf{lm}=\lim b_k.$
The dichotomic algorithm  is implemented in the (matlab) mex-file median\_rouen\_dic given in the supplementary material.
Table~1 contains  computational times of median\_rouen\_dic and those of the build-in matlab function median
when the sample is ``small". This table shows clearly that, the performance of  median\_rouen\_dic is superior
to that of matlab median when the sample size is between 10 and 150.

From Table~2 we see that, the runtime of the program median\_rouen\_dic increases with the sample size $n.$ This is because the number of 
comparisons increase with $n.$  The runtime of matlab median 
also increases with the sample size. It also appears from this table that the performance of median\_rouen\_dic is slightly
superior
to that of matlab median for ``medium" sample sizes.
\begin{table}[!htbp]
	\caption{Runtimes in seconds of  median\_rouen\_dic and the build-in matlab median  over $500$ replications.}
	\label{table1}
	\begin{tabular}{lccc}
		\toprule
		\multicolumn{4}{c}{\textsc{Uniform over $(0,\ 1)$\quad\textit{small sample sizes}} }\\
		\toprule
		{} & \specialcell{Time of\\ median\_rouen\_dic} & \specialcell{Time of \\ matlab median} & Ratio\\
		\midrule
		n=10 & 0.00278   & 0.00287  & 0.9676 \\
		\midrule
		n=30 & 0.00323   & 0.00343  & 0.9417 \\
		\midrule
		n=50 & 0.00495   & 0.00539  & 0.9170 \\
		\midrule
		n=70 & 0.00297   & 0.00330  & 0.8986 \\
		\midrule
		n=90 & 0.00364   & 0.00399  & 0.9137 \\
		\midrule
		n=110 & 0.00494  & 0.00543  & 0.9099\\
		\midrule
		n=130 & 0.00473  & 0.00527  & 0.8975 \\
		\midrule
		n=150 & 0.00431  & 0.00482  & 0.8949 \\
		\bottomrule
	\end{tabular}
\end{table}


\begin{table}[!htbp]
	\caption{Runtimes in seconds of  median\_rouen\_dic and the build-in matlab median  over $500$ replications.}
	\label{table2}
	\begin{tabular}{lccc}
		\toprule
		\multicolumn{4}{c}{\textsc{Uniform over $(0,\ 1)$\quad\textit{Medium sample sizes}} }\\
		\toprule
		{} & \specialcell{Time of\\ median\_rouen\_dic} & \specialcell{Time of \\ matlab median} & Ratio\\
		\midrule
		n=5000 & 0.04232  & 0.04650    &  0.9102       \\
		\midrule
		n=8000 & 0.06728  & 0.07394    &  0.9099      \\
		\midrule
		n=11000 & 0.08623 & 0.09419    &  0.9155    \\
		\midrule
		n=14000 & 0.10798 & 0.11757    &  0.9184  \\
		\midrule
		n=17000 & 0.12838  & 0.14049   &  0.9138  \\
		\midrule
		n=20000 & 0.15256  & 0.16782   & 0.9090     \\
		\midrule
		n=23000 & 0.17211  & 0.18817   & 0.9147     \\
		\midrule
		n=26000 & 0.19742  & 0.21676   & 0.9107     \\
		\bottomrule
	\end{tabular}
\end{table}





\begin{table}[!htbp]
	\caption{Runtimes in seconds of  median\_rouen\_dic and the build-in matlab median  over $500$ replications.}
	\label{table3}
	\begin{tabular}{lccc}
		\toprule
		\multicolumn{4}{c}{\textsc{Uniform over $(0,\ 1)$\quad\textit{large sample sizes}} }\\
		\toprule
		{} & \specialcell{Time of\\ median\_rouen\_dic} & \specialcell{Time of \\ matlab median} & Ratio\\
		\midrule
		n=1000000 & 10.2957 & 8.8177   &  1.1676   \\
		\midrule
		n=2000000 & 21.4221 & 22.3411  &  0.9589 \\
		\midrule
		n=3000000  & 31.9876 & 33.2991 &  0.9606 \\
		\midrule
		n=4000000 &  36.9473 & 36.4557   & 1.0135 \\
		\midrule
		n=5000000  & 62.2457 &  67.5856  & 0.9210  \\
		\midrule
		n=6000000  & 75.2828   & 81.3021 & 0.9260   \\
		\midrule
		n=7000000 &  85.8764   & 92.7082 & 0.9263   \\
		\midrule
		n=8000000 &  98.4278  &  106.6332 &  0.9230    \\
		\bottomrule
	\end{tabular}
\end{table}

In Table ~3 is displayed the results of the computations when the sample sizes are ``large". As in Table~2, the runtimes of
both methods increase with the sample size. In this table, no method   is superior to the other. However
median\_rouen\_dic  is superior when the
sample size is between 5 000 000 and 8 000 000.



\section{Gradient descent}\label{sec:4}








\subsection{Plain descent method }\label{sec:4.1}

This approach borrows a lot from the classical gradient descent algorithm. Before its introduction, let us
recall the basic tools of the gradient descent. Let $f:\R\longrightarrow \R$ be a convex
differentiable function. Assume that we want to find a minimizer  of $f.$ The gradient descent
method consists in building a sequence $(x_n)_{n\in\N}$ converging to a  minimizer of $f.$
The sequence is constructed as follows:
\begin{itemize}
	\item Select an arbitrary starting point $x_0;$
	\item Choose a learning rate $\delta >0;$
	\item The sequence is defined recursively as
	\[x_{n+1}=x_n -\delta f'(x_n).\]
\end{itemize}
%%%%%%%%%
The theory shows that, under certain regularity conditions of $f,$ there exists $\delta_0>0$ 
such that
\[
0<\delta<\delta_0 \implies x_n\quad\text{converges to}\quad w_0
\]
where $w_0$ is a minimizer of $f.$ Some comments regarding the construction of the sequence
are in order here. If $f'(x_n)>0,$ then f is increasing in a neighborhood of $x_n.$ Thus, the minimizer
of $f$ is less than $x_n,$ this is why $x_{n+1}$ is chosen such that
\[
x_{n+1}=x_n -\delta f'(x_n)<x_n.
\]
If $f'(x_n)<0,$ then
\[
x_n<  x_{n+1}=x_n -\delta f'(x_n).
\]
This means that $ x_{n+1}$ is chosen to the left of $x_n$ when the gradient is positive and to
the right when the gradient is negative: we move against the gradient. We now illustrate the working
of the gradient descent on a simple example. Let $a$ be a fixed real and $f(x)={\frac{1}{2}(x-a)^2}.$
It is clear that $f$ has a unique minimum at $w_0=a.$ We have:\quad $f'(x)=x-a,$ therefore
the sequence $(x_n)_{n\in \N}$ is defined recursively by
\begin{align*} 
	x_{n+1} & =x_n-\delta (x_n-a) \\
	& =(1-\delta) x_n+ \delta a.      
\end{align*}
Because $a=(1-\delta) a+\delta a,$ we see that
\[
x_{n+1}-a   =(1-\delta) (x_n-a)\]
it follows that
\[
x_n-a=(1-\delta)^n (x_0-a)\iff x_n=a+ (1-\delta)^n (x_0-a).
\]        
Therefore we have:
\[
\lim x_n=a \iff -1<1-\delta <1\iff 0<\delta<2.
\]


We introduce the following definition.    
\begin{definition}
	A sequence $\left(\mathscr{B}_k=\left(a_k,b_k,c_k,M_k \right),\quad k\in \N^*\right)$ is called a gradient descent 
	sequence if it is built recursively by  Step-b1  and  Step-b2 given in the supplementary material.
\end{definition}    



Proposition 5 below will allow us to state some properties of a gradient descent sequence.
Let $a<b,$ and set $m=\mathrm{card}(\{i:\ a\leq X_i\leq b\})$. We have the following useful result.

\begin{proposition} If $\varphi_{-}'(a)<0$ and $\left(\varphi_{-}'(b)>0\right.$ or 
	$\left( \varphi_{-}'(b)=0 \right.$ and $\left.\left.\varphi_{+}'(b)>0\right)\right)$ then
	%%%%%%%% 
	\[
	-1<\frac{\varphi_{-}'(a)}{2m}<0.
	\]
\end{proposition}
The proof of Proposition 5 is in the supplementary material.


We now present the intuition behind the gradient descent sequence. Let $a<b,$ we assume that $\varphi_{-}'(a)<0 \quad\text{and}\quad \varphi_{+}'(b)>0$,
so that $\mathbf{lm}, \mathbf{gm}\in [a,\ b].$  We need the following notations
\[
M=(a\leq M_1\leq b)\quad\text{where}\quad M_1=(X_1,\ldots,X_n)
\] 
thus $X_i$ is a component of the vector M if $a\leq X_i\leq b$. To save space we set
\[
M=(Z_1,\ldots,Z_m)
\]
therefore we assume that $\mathtt{length}(M)=m=\mathrm{card}\{i:\ a\leq X_i\leq b\}.$ In what follows it is
important to remember that $\mathbf{lm}$ and $\mathbf{gm}$ are components of the vector $M.$
Given $a,b,M,m$ defined above, the gradient descent
algorithm chooses $c\in(a,\ b)$ by the formula
\[
c= a -\frac{\varphi_{-}'(a)}{2m} (b-a).
\]
We want to explain how this formula is obtained. We also assume that $n=2p$ and $\mathbf{lm}<\mathbf{gm},$ because
the intuition is easier to ``see'' in this situation. We set
\begin{align*}
	r_M(\mathbf{lm})= & \sum_{j=1}^m\indica_{[Z_i \leq \mathbf{lm} ]},\\
	t_0=a, & \quad t_j=a+ \frac{j}{m}(b-a),\quad j=1,\ldots,m.
\end{align*}
The gradient descent algorithm chooses $t_{j_0}$ in $(a,\ b)$ if
\[
r_M(\mathbf{lm})=j_0.
\]
This means that the chosen point is ``close'' to $b$ when $r_M(\mathbf{lm})$ is large and ``close'' to $a$
when $r_M(\mathbf{lm})$ is small. Indeed,  let $c\in (\mathbf{lm},\  \mathbf{gm}),$ we have
\begin{align*}
	\varphi_{-}'(c)=0 &=\varphi_{-}'(a)+2\sum_{i=1}^n \indica_{[X_i <c]} \indica_{[a\leq X_i \leq b]}\\
	&=\varphi_{-}'(a)+2\sum_{j=1}^m \indica_{[Z_i <c]}\\ 
	&=\varphi_{-}'(a)+2 r_M(\mathbf{lm}). 
\end{align*}
It follows that
\[
r_M(\mathbf{lm})=-\frac{\varphi_{-}'(a)}{2}
\]
hence
\begin{align*}
	c= & a -\frac{\varphi_{-}'(a)}{2m} (b-a)\\
	= & a+\frac{r_M(\mathbf{lm})}{m}(b-a)\\
	= & t_{r_M(\mathbf{lm})}.
\end{align*}
Why do we talk about gradient descent (regarding this approach)? Because when $a_k<b_k,$ $a_{k+1},b_{k+1}$ and $M_{k+1}$
are derived from $\left(a_k,b_k,c_k,M_k \right)$ using $\varphi_{-}'(c_k).$ Moreover by setting
\begin{align*}
	L_k=& (M_k<c_k),\\
	R_k= &(M_k\geq c_k),\\
	m_k=&\mathtt{length}(M_k),\\
	\lambda_k=&(b_k-a_k)/2m_k,
\end{align*}
if $(\varphi_{-}'(c_k)<0)$ and $(L_k\neq \emptyset)$ we have:
\begin{align*}
	a_{k+1}=& c_k\\
	= & a_k-\lambda_k \varphi_{-}'(a_k).
\end{align*}
And, this formula resembles the classical gradient descent formula.\par\medspace




We  describe below two situations in which a gradient descent sequence converges to \textbf{lm}.

\noindent{\textsc{--Situation 1:}} If $n=2p+1$\par

By Corollary 1 we know that:
\[
\mathbf{lm}=\mathbf{am},\quad \varphi_{-}'(\mathbf{lm})<0,\quad \varphi_{+}'(\mathbf{lm})>0.
\]
Since by construction $a_k\leq \mathbf{lm}\leq b_k,$ we deduce that
\[
\varphi_{-}'(a_k)<0 \quad\text{and}\quad \varphi_{+}'(b_k)>0,\quad \forall k\in \N^*,
\]
because $x\mapsto \varphi_{-}'(x)$ and $x\mapsto \varphi_{+}'(x)$ are nondecreasing functions.
If $a_k<b_k,$ then by the gradient descent algorithm we have:
\begin{align*}
	c_k=& a_k-r\rho(a_k)\\
	=& a_k-r\varphi_{-}'(a_k)\quad (\text{definition of}\ \rho )\\
	=& a_k -\frac{\varphi_{-}'(a_k)}{2m} (b_k-a_k).\quad  (m=\mathtt{length}( M_k))
\end{align*}
Because of Proposition 5 we have 
$0< -\frac{\varphi_{-}'(a_k)}{2m}<1,$ thus we deduce that $c_k \in (a_k,\ b_k).$  This fact shows that the gradient descent sequence
is a generic sequence. Therefore the sequences $\left(a_k\right)_{k\in \N^*}$ and $\left(b_k\right)_{k\in \N^*}$ converge
to \textbf{lm.}

\noindent{\textsc{--Situation 2:}} If $n=2p$ and $\mathbf{lm}=\mathbf{gm}$\par 
In this case we know that $\mathbf{lm}=\mathbf{am},$ so by Theorem 1 we have:
\[
\varphi_{-}'(\mathbf{lm})<0,\quad \varphi_{+}'(\mathbf{lm})>0.
\]
Therefore, as in the case $n=2p+1$ we have
\[
\lim a_k=\mathbf{lm}=\lim b_k.
\]  

The case $n=2p$ and $\mathbf{lm}<\mathbf{gm} $ is more difficult  to analyze. When this condition is true,
we know that (see the proof of Theorem 1) 
\begin{align*}
	\varphi_{-}'(\mathbf{lm})< & 0,\quad \varphi_{+}'(\mathbf{lm})=0,\\
	\varphi'(t)= &0,\quad \forall t\in (\mathbf{lm},\ \mathbf{gm}),\\
	\varphi_{-}'(\mathbf{gm})= & 0,\quad \varphi_{+}'(\mathbf{gm})>0.
\end{align*}
Using inequalities and equalities above, we can prove that
\begin{align*}
	\varphi_{-}'(t)< & 0,\quad \text{for}\quad t\in (-\infty,\ \mathbf{lm}),\quad  
	\varphi_{-}'(t)= 0,\quad \text{for}\quad t\in (\mathbf{lm},\ \mathbf{gm}],\quad\\
	\varphi_{-}'(t)> & 0,\quad \text{for}\quad t\in (\mathbf{gm},\ +\infty).
\end{align*}

The set $W_0$ defined by 
\[
W_0=(\mathbf{lm},\ \mathbf{gm}]
\] 
and Definition 4 introduced below play important roles in the rest of this subsection.            
\begin{definition}
	A sequence $\left(\mathscr{C}_k=\left(a_k,b_k,c_k,M_k \right),\quad k\in \N^*\right)$ is called a mix	sequence if
	it is built recursively by  Step-c1  and  Step-c2 presented in the supplementary material.
\end{definition}    



Lemma 5--7 below will help us to prove some properties of the gradient descent sequence. Their proofs are in the supplementary material.

\begin{lemma} Let $\left(\mathscr{C}_k=\left(a_k,b_k,c_k,M_k \right),\quad k\in \N^*\right)$ be a mix sequence.
	If $n=2p$ and $\mathbf{lm}<\mathbf{gm},$ we have:
	\[
	c_1,\ldots,c_q\notin W_0\implies a_1\leq\cdots \leq a_{q+1}\leq \mathbf{lm}<\mathbf{gm}\leq b_{q+1}\leq \cdots \leq b_1.
	\]
	
\end{lemma}	



\begin{lemma}  Let $\left(\mathscr{C}_k=\left(a_k,b_k,c_k,M_k \right),\quad k\in \N^*\right)$ be a mix sequence.
	If $n=2p$ and $\mathbf{lm}<\mathbf{gm},$ we have:
	\[
	a_k<b_k\implies c_k\in (a_k,\ b_k).
	\]
\end{lemma}



\begin{lemma}  Let $\left(\mathscr{C}_k=\left(a_k,b_k,c_k,M_k \right),\quad k\in \N^*\right)$ be a mix sequence.
	If $n=2p$ and $\mathbf{lm}<\mathbf{gm},$ there exists $k_0\in \N^*$ such that:
	\[
	c_{k_0}\in (\mathbf{lm},\ \mathbf{gm}].
	\]
\end{lemma}

The property of the mix sequence stated in Lemma 7 is also valid for the gradient descent sequence. This Property is given
in Proposition 6 below. The proof of Proposition 6 is in the supplementary material.
\begin{proposition}  Let $\left(\mathscr{B}_k=\left(a_k,b_k,c_k,M_k \right),\quad k\in \N^*\right)$ be a gradient descent sequence.
	If $n=2p$ and $\mathbf{lm}<\mathbf{gm},$ there exists $q_0\in \N^*$ such that:
	\[
	c_{q_0}\in (\mathbf{lm},\ \mathbf{gm}].
	\]
\end{proposition}




Proposition 6 is very important for the gradient descent approach when $n=2p$ and $\mathbf{lm}<\mathbf{gm}.$
Let $\left(\mathscr{B}_k=\left(a_k,b_k,c_k,M_k \right),\quad k\in \N^*\right)$ be a gradient descent sequence and
set 
\[
q_0=\min\{ k,\quad c_k\in (\mathbf{lm},\ \mathbf{gm}]\}.
\]
In the implementation of the gradient descent we will only build
\[
\left(a_1,b_1,c_1,M_1 \right),\dots,\left(a_{q_0},b_{q_0},c_{q_0},M_{q_0} \right).
\]
Because we know how to find $\mathbf{lm}$ and  $\mathbf{gm}$ in $M_{q_0}$ when $c_{q_0}$ is found.



The gradient descent algorithm is implemented in the program median\_rouen\_gra given in the supplementary material. In Tables 4-6 are presented the
results of our experimentations  designed to compare matlab median and median\_rouen\_gra. The performance of
median\_rouen\_gra (relatively to that of matlab median) is grossly similar to that of median\_rouen\_dic.
Table 4 shows that the runtime of median\_rouen\_gra is better than that of matlab median when the sample size is ``small";
Table 5 shows superiority of median\_rouen\_gra to matlab median for ``medium" sample sizes; in Table 6 matlab median
is superior to median\_rouen\_gra when the sample size is 1 million  while
median\_rouen\_gra performs better than matlab median when the sample size is between 2 million and 8 million.







\begin{table}[!htbp]
	\caption{Runtimes in seconds of  median\_rouen\_gra and the build-in matlab median  over $500$ replications.}
	\label{table4}
	\begin{tabular}{lccc}
		\toprule
		\multicolumn{4}{c}{\textsc{Uniform over $(0,\ 1)$\quad\textit{small sample sizes}} }\\
		\toprule
		{} & \specialcell{Time of\\ median\_rouen\_gra} & \specialcell{Time of \\ matlab median} & Ratio\\
		\midrule
		n=10 & 0.00242  & 0.00268   & 0.9009  \\
		\midrule
		n=30 & 0.00272  & 0.00322   & 0.8460 \\
		\midrule
		n=50 & 0.00345  & 0.00416   & 0.8288  \\
		\midrule
		n=70 & 0.00499  & 0.00692   & 0.7214   \\
		\midrule
		n=90 & 0.00435  & 0.00538   & 0.8085  \\
		\midrule
		n=110 & 0.00374 & 0.00467   & 0.8007   \\
		\midrule
		n=130 & 0.00287 & 0.00361   & 0.7983   \\
		\midrule
		n=150 & 0.00411 & 0.00525   & 0.7823    \\
		\bottomrule
	\end{tabular}
\end{table}


\begin{table}[!htbp]
	\caption{Runtimes in seconds of  median\_rouen\_gra and the build-in matlab median  over $500$ replications.}
	\label{table5}
	\begin{tabular}{lccc}
		\toprule
		\multicolumn{4}{c}{\textsc{Uniform over $(0,\ 1)$\quad\textit{Medium sample sizes}} }\\
		\toprule
		{} & \specialcell{Time of\\ median\_rouen\_gra} & \specialcell{Time of \\ matlab median} & Ratio\\
		\midrule
		n=5000  & 0.03214  & 0.0434  & 0.7396   \\
		\midrule
		n=8000  & 0.0496   & 0.0680  & 0.7300  \\
		\midrule
		n=11000 & 0.06773  & 0.09323 & 0.7265 \\
		\midrule
		n=14000 & 0.08510  & 0.11709 & 0.7267   \\
		\midrule
		n=17000 & 0.10142  & 0.13994 & 0.7248  \\
		\midrule
		n=20000 & 0.11921  & 0.16434 & 0.7254  \\
		\midrule
		n=23000 & 0.13600  & 0.18849 & 0.7215  \\
		\midrule
		n=26000 & 0.15546  & 0.21637 & 0.7185   \\
		\bottomrule
	\end{tabular}
\end{table}





\begin{table}[!htbp]
	\caption{Runtimes in seconds of  median\_rouen\_grad and the build-in matlab median  over $500$ replications.}
	\label{table6}
	\begin{tabular}{lccc}
		\toprule
		\multicolumn{4}{c}{\textsc{Uniform over $(0,\ 1)$\quad\textit{large sample sizes}} }\\
		\toprule
		{} & \specialcell{Time of\\ median\_rouen\_gra} & \specialcell{Time of \\ matlab median} & Ratio\\
		\midrule
		n=1000000 &  9.3371  &  8.8944  &  1.0498  \\
		\midrule
		n=2000000 &  16.9671 &  18.9463 &  0.8955 \\
		\midrule
		n=3000000 & 25.3496  & 28.5725 &  0.8872  \\
		\midrule
		n=4000000 & 39.4218  & 45.5724 &  0.8650   \\
		\midrule
		n=5000000 & 56.9510  & 67.5012 &  0.8437   \\
		\midrule
		n=6000000 & 67.8968  & 79.9994 &  0.8487  \\
		\midrule
		n=7000000 & 78.2593  & 93.1864 &  0.8398   \\
		\midrule
		n=8000000 & 90.1274  & 106.5925 & 0.8455   \\
		\bottomrule
	\end{tabular}
\end{table}

A comparison of the third columns of Tables 1-3 and those of Tables 4-5 reveals that the performance of median\_rouen\_gra
is superior to that of median\_rouen\_dic. 



\subsection{Truncated gradient descent }\label{sec:4.2}

This approach is a modification of the gradient descent method. In this subsection, we only mention how the gradient descent is modified.
Let $\left(\mathscr{G}_k=\left(a_k,b_k,c_k,M_k \right),\quad k\in \N^*\right)$ be the gradient descent sequence, let us
recall how $c_k$ is ``chosen" in $(a_k,\ b_k)$ when  $ \varphi_{-}'(a_k)<0\quad\text{and}\quad \varphi_{+}'(b_k)>0$: $c_k$ is given
by 
\[
c_k=a_k+h_k (b_k-a_k)\quad\text{where}\quad h_k=-\left(\frac{\varphi_{-}'(a_k)}{2m_k}\right),\quad
m_k=\mathtt{length}(M_k).
\]
Let us also recall the following vectors useful in the sequel:
\[
R=R_k=\left(M_k\geq c_k\right),\quad  L=L_k=\left(M_k< c_k\right).
\]
It follows from the formula of $c_k$ that, if $h_k$ is close to 1, then $c_k$ is close $b_k.$ When $c_k$ is close
to $b_k$ the likelihood of getting $ R=R_k=\emptyset$ is high. And, when $R=\emptyset$ we know that
\[
b_{k+1}=\max(L)
\]
thus the computation of $b_{k+1}$  is time consuming. If $h_k$ is close to 0, the risk of getting $ L=L_k=\emptyset$
is also high. When $L=\emptyset$ the computation of $a_{k+1}$ is costly. To avoid these situations from happening,
we truncate $h_k$ in the following manner. Let $\alpha\in(0,\ \frac{1}{2})$ we set
\[
\bar{h}_k = \begin{cases} \alpha, & \text{if}\quad h_k < \alpha, \\ 
	h_k, & \text{if}\quad\alpha \leq h_k\leq 1-\alpha, \\
	1-\alpha, &\text{if}\quad h_k > 1-\alpha, 
\end{cases} 
\]
and, $c_k$ in the truncated gradient descent is defined by
\[
c_k=a_k+\bar{h}_k (b_k-a_k).
\]


\begin{table}[!htbp]
	\caption{Runtimes in seconds of  median\_rouen\_tru and the build-in matlab median  over $500$ replications.}
	\label{table7}
	\begin{tabular}{lccc}
		\toprule
		\multicolumn{4}{c}{\textsc{Uniform over $(0,\ 1)$\quad\textit{small sample sizes}} }\\
		\toprule
		{} & \specialcell{Time of\\ median\_rouen\_tru} & \specialcell{Time of \\ matlab median} & Ratio\\
		\midrule
		n=10 & 0.00273 &  0.00303 &   0.9004   \\
		\midrule
		n=30 & 0.00409 &  0.00452 &   0.9047   \\
		\midrule
		n=50 & 0.00451 &  0.00508 &   0.8863    \\
		\midrule
		n=70 & 0.0035  &  0.00402 &   0.8829    \\
		\midrule
		n=90 & 0.00332 &  0.00374 &   0.8867    \\
		\midrule
		n=110 & 0.00424 &  0.00477 &  0.8883  \\
		\midrule
		n=130 & 0.00388 &  0.00439 &  0.8838   \\
		\midrule
		n=150 & 0.00345 &  0.00391 &  0.8832     \\
		\bottomrule
	\end{tabular}
\end{table}


\begin{table}[!htbp]
	\caption{Runtimes in seconds of  median\_rouen\_tru and the build-in matlab median  over $500$ replications.}
	\label{table8}
	\begin{tabular}{lccc}
		\toprule
		\multicolumn{4}{c}{\textsc{Uniform over $(0,\ 1)$\quad\textit{Medium sample sizes}} }\\
		\toprule
		{} & \specialcell{Time of\\ median\_rouen\_tru} & \specialcell{Time of \\ matlab median} & Ratio\\
		\midrule
		n=5000 &  0.04099 &  0.04355  & 0.9412  \\
		\midrule
		n=8000 &  0.06385 &  0.06744  & 0.9468   \\
		\midrule
		n=11000 & 0.08767 &  0.09318  & 0.9409   \\
		\midrule
		n=14000&  0.10894 &  0.11599  & 0.9392  \\
		\midrule
		n=17000 & 0.13220 &  0.14026  & 0.9425  \\
		\midrule
		n=20000 & 0.15422 &  0.16433  & 0.9385   \\
		\midrule
		n=23000 & 0.17664 &  0.18831  & 0.9380  \\
		\midrule
		n=26000 & 0.19900 &  0.21225  & 0.9375   \\
		\bottomrule
	\end{tabular}
\end{table}





\begin{table}[!htbp]
	\caption{Runtimes in seconds of  median\_rouen\_tru and the build-in matlab median  over $500$ replications.}
	\label{table9}
	\begin{tabular}{lccc}
		\toprule
		\multicolumn{4}{c}{\textsc{Uniform over $(0,\ 1)$\quad\textit{large sample sizes}} }\\
		\toprule
		{} & \specialcell{Time of\\ median\_rouen\_tru} & \specialcell{Time of \\ matlab median} & Ratio\\
		\midrule
		n=1000000 & 7.60236   &   8.61220  & 0.8827  \\
		\midrule
		n=2000000  & 15.61259 &  17.92161  & 0.8711 \\
		\midrule
		n=3000000  & 23.21215 &  27.04747  & 0.8582   \\
		\midrule
		n=4000000 &  32.82092 &  38.50883  & 0.8523   \\
		\midrule
		n=5000000  & 50.79466 &  62.04248  & 0.8187  \\
		\midrule
		n=6000000  & 68.00246 &  86.36158  & 0.7874   \\
		\midrule
		n=7000000 &  73.83915 & 90.48521   & 0.8160 \\
		\midrule
		n=8000000 &  82.72599 & 101.02673  & 0.8188  \\
		\bottomrule
	\end{tabular}
\end{table}

The truncated gradient descent algorithm is implemented in the matlab (mex) program median\_rouen\_tru presented
in the supplementary material. In Tables 7--9 are summarized our computations regarding the truncated gradient descent method.
These tables reveal that the runtime of  median\_rouen\_tru is better than that of the built-in matlab median for
all samples sizes.

The truncation parameter $\alpha$ used in the truncated method is supposed to satisfy $0<\alpha<\frac{1}{2}.$ 
If we ``choose" $\alpha=\frac{1}{2}$ we get the dichotomic approach, if $\alpha=0$ we get the ``plain gradient" algorithm.
This observation shows that each one of our approaches can be considered a special case of the truncated gradient method.
When the samples sizes are ``small'' the best results are obtained by plain gradient  and the  
truncated gradient. The best results
are obtained by the plain gradient  approach and dichotomy when the samples sizes are ``medium'', and the ``truncated gradient" has 
the best results when $n$ is ``large." Those facts  suggest that we should ``choose'' $\alpha$  as a function of the sample size.
We hope to consider the problem of selecting $\alpha=\alpha_n$ in a future paper.







Assume that the medians compute by the matlab build-in function are stored in the vector
$V_{mat}=(u_1,\ldots,u_r),$ where $r$ is the number of replications. Similarly suppose that
the medians computed by median\_rouen\_dic (respectively median\_rouen\_grad, median\_rouen\_tru) are stored in the
vector $V_{rou}=(v_1,\ldots,v_r).$ In the simulations we have also computed the quantities:
\[
Error=\sum_{i=1}^{r} (u_i-v_i)^2.
\]
These quantities are not reported because they are all equal to zero. This shows that the methods introduced in this paper
give the same results as the matlab build-in function.

The authors report there are no competing interests to declare.

Data sharing is not applicable to this article. The data used in this
paper are RANDOM samples generated by Matlab software, thus can be
created by any reader.

\protect\phantomsection\label{refs}
\begin{CSLReferences}{1}{1}
\bibitem[\citeproctext]{ref-AzPeTo:23}
Azzini, Ivano, Domenico Perrotta, and Francesca Torti. 2023. \emph{A
Practically Efficient Fixed-Pivot Selection Algorithm and Its Extensible
MATLAB Suite}. \url{https://arxiv.org/abs/2302.05705}.

\bibitem[\citeproctext]{ref-Hoa:61}
Hoare, C. A. R. 1961. {``Algoritm 65: Find.''} \emph{Communications of
the ACM} 4 (7): 321--22.

\bibitem[\citeproctext]{ref-HuYaTa:79}
Huang, T., G. Yang, and G. Tang. 1979. {``A Fast Two-Dimensional Median
Filtering Algorithm.''} \emph{Acoustics, Speech and Signal Processing,
IEEE Transactions on} 27 (1): 13--18.
\url{https://doi.org/10.1109/TASSP.1979.1163188}.

\bibitem[\citeproctext]{ref-Ma:19}
Ma, Wing-Kin (Ken). 2019. \emph{Subgradient Methods}.
\url{https://arxiv.org/abs/0806.3301}.

\bibitem[\citeproctext]{ref-NoGa:82}
Nodes, T., and N. Gallagher. 1982. {``Median Filters: Some Modifications
and Their Properties.''} \emph{IEEE Transactions on Acoustics, Speech,
and Signal Processing} 30 (5): 739--46.
\url{https://doi.org/10.1109/TASSP.1982.1163951}.

\bibitem[\citeproctext]{ref-Sat:22}
Satman, Mehmet Hakan. 2022. {``{Teaching the Median with Terms of
Absolute Value, Differentiability, and Optimization}.''}
\emph{Alphanumeric Journal} 10 (1): 41--50.
https://doi.org/\url{https://doi.org/10.17093/}.

\bibitem[\citeproctext]{ref-Ste:85}
Stewart, Robert R. 1985. {``{Median Filtering: Review and a new f/k
analogue design }.''} \emph{Journal of the Canadian Society of
Exploration Geophysicists} 21 (1): 54--63.

\bibitem[\citeproctext]{ref-Tib:09}
Tibshirani, Ryan J. 2009. \emph{Fast Computation of the Median by
Successive Binning}. \url{https://arxiv.org/abs/0806.3301}.

\bibitem[\citeproctext]{ref-Tib:17}
Tibshirani, Ryan J. 2017. \emph{Subgradient Method, Convex Optimization
10-725}. \url{https://arxiv.org/abs/0806.3301}.

\bibitem[\citeproctext]{ref-YINeALL:96}
Yin, Lin, Ruikang Yang, M. Gabbouj, and Y. Neuvo. 1996. {``Weighted
Median Filters: A Tutorial.''} \emph{IEEE Transactions on Circuits and
Systems II: Analog and Digital Signal Processing} 43 (3): 157--92.
\url{https://doi.org/10.1109/82.486465}.

\end{CSLReferences}




\end{document}
